% problem_id: S001-lemma-algebraic-quasi-projective
% source_id: S001 (Stacks Project)
% source_file: groupoids.tex
% statement_locator: groupoids.tex lines 1364--1368
% proof_locator: groupoids.tex lines 1370--1480
% env: lemma
% id_note: label 在本来源内唯一
% pairing: tight (verified)
% status: complete (statement + author proof verbatim + reason + locators)
%
\begin{lemma}
\label{lemma-algebraic-quasi-projective}
Let $k$ be a field. Let $G$ be an algebraic group scheme over $k$.
Then $G$ is quasi-projective over $k$.
\end{lemma}

\begin{proof}
By Varieties, Lemma \ref{varieties-lemma-ample-after-field-extension}
we may assume that $k$ is algebraically closed. Let $G^0 \subset G$
be the connected component of $G$ as in
Proposition \ref{proposition-connected-component}.
Then every other connected component of $G$ has a $k$-rational
point and hence is isomorphic to $G^0$ as a scheme.
Since $G$ is quasi-compact and Noetherian, there are finitely many of these
connected components. Thus we reduce to the case discussed in
the next paragraph.

\medskip\noindent
Let $G$ be a connected algebraic group scheme over an algebraically closed
field $k$. If the characteristic of $k$ is zero, then $G$ is smooth over
$k$ by Lemma \ref{lemma-group-scheme-characteristic-zero-smooth}.
If the characteristic of $k$ is $p > 0$, then we let $H = G_{red}$
be the reduction of $G$. By
Divisors, Proposition \ref{divisors-proposition-push-down-ample}
it suffices to show that $H$ has an ample invertible sheaf.
(For an algebraic scheme over $k$ having an ample invertible
sheaf is equivalent to being quasi-projective over $k$, see
for example the very general
More on Morphisms, Lemma \ref{more-morphisms-lemma-quasi-projective}.)
By Lemma \ref{lemma-reduced-subgroup-scheme-perfect}
we see that $H$ is a group scheme over $k$.
By Lemma \ref{lemma-reduced-group-scheme-prefect-field-characteristic-p-smooth}
we see that $H$ is smooth over $k$.
This reduces us to the situation discussed in the next
paragraph.

\medskip\noindent
Let $G$ be a quasi-compact irreducible smooth group scheme over an
algebraically closed field $k$. Observe that the local rings of $G$
are regular and hence UFDs
(Varieties, Lemma \ref{varieties-lemma-smooth-regular} and
More on Algebra, Lemma \ref{more-algebra-lemma-regular-local-UFD}).
The complement of a nonempty affine open of $G$
is the support of an effective Cartier divisor $D$.
This follows from Divisors, Lemma
\ref{divisors-lemma-complement-open-affine-effective-cartier-divisor}.
(Observe that $G$ is separated by
Lemma \ref{lemma-group-scheme-over-field-separated}.)
We conclude there exists an effective Cartier divisor $D \subset G$
such that $G \setminus D$ is affine. We will use below that
for any $n \geq 1$ and $g_1, \ldots, g_n \in G(k)$ the complement
$G \setminus \bigcup D g_i$ is affine. Namely, it is the intersection
of the affine opens $G \setminus Dg_i \cong G \setminus D$
in the separated scheme $G$.

\medskip\noindent
We may choose the top row of the diagram
$$
\xymatrix{
G & U \ar[l]_j \ar[r]^\pi & \mathbf{A}^d_k \\
& W \ar[r]^{\pi'} \ar[u] & V \ar[u]
}
$$
such that $U \not = \emptyset$, $j : U \to G$ is an open immersion, and
$\pi$ is \'etale, see
Morphisms, Lemma \ref{morphisms-lemma-smooth-etale-over-affine-space}.
There is a nonempty affine open $V \subset \mathbf{A}^d_k$ such that
with $W = \pi^{-1}(V)$ the morphism $\pi' = \pi|_W : W \to V$ is finite \'etale.
In particular $\pi'$ is finite locally free, say of degree $n$.
Consider the effective Cartier divisor
$$
\mathcal{D} = \{(g, w) \mid m(g, j(w)) \in D\} \subset G \times W
$$
(This is the restriction to $G \times W$ of the pullback of $D \subset G$
under the flat morphism $m : G \times G \to G$.)
Consider the closed subset\footnote{Using the material
in Divisors, Section \ref{divisors-section-norms}
we could take as effective Cartier
divisor $E$ the norm of the effective Cartier divisor $\mathcal{D}$
along the finite locally free morphism $1 \times \pi'$ bypassing
some of the arguments.}
$T = (1 \times \pi')(\mathcal{D}) \subset G \times V$.
Since $\pi'$ is finite locally free, every irreducible component
of $T$ has codimension $1$ in $G \times V$. Since $G \times V$
is smooth over $k$ we conclude these components are effective Cartier
divisors (Divisors, Lemma \ref{divisors-lemma-weil-divisor-is-cartier-UFD}
and lemmas cited above)
and hence $T$ is the support of an effective Cartier divisor
$E$ in $G \times V$. If $v \in V(k)$, then
$(\pi')^{-1}(v) = \{w_1, \ldots, w_n\} \subset W(k)$ and we see that
$$
E_v = \bigcup\nolimits_{i = 1, \ldots, n} D j(w_i)^{-1}
$$
in $G$ set theoretically. In particular we see that $G \setminus E_v$
is affine open (see above).
Moreover, if $g \in G(k)$, then there exists a $v \in V$
such that $g \not \in E_v$. Namely, the set $W'$ of $w \in W$ such that
$g \not \in Dj(w)^{-1}$ is nonempty open and it suffices to pick $v$
such that the fibre of $W' \to V$ over $v$ has $n$ elements.

\medskip\noindent
Consider the invertible sheaf
$\mathcal{M} = \mathcal{O}_{G \times V}(E)$ on $G \times V$.
By Varieties, Lemma \ref{varieties-lemma-rational-equivalence-for-Pic}
the isomorphism class $\mathcal{L}$ of the restriction
$\mathcal{M}_v = \mathcal{O}_G(E_v)$ is independent of $v \in V(k)$.
On the other hand, for every $g \in G(k)$ we can find a $v$
such that $g \not \in E_v$ and such that $G \setminus E_v$
is affine. Thus the canonical section
(Divisors, Definition
\ref{divisors-definition-invertible-sheaf-effective-Cartier-divisor})
of $\mathcal{O}_G(E_v)$
corresponds to a section $s_v$ of $\mathcal{L}$ which does not
vanish at $g$ and such that $G_{s_v}$ is affine.
This means that $\mathcal{L}$ is ample by definition
(Properties, Definition \ref{properties-definition-ample}).
\end{proof}
